\documentclass[10pt,a4paper]{amsart}
\usepackage{amsmath,amssymb,amsthm,graphicx}
\usepackage{verbatim}
\newenvironment{reference}{\comment}{\endcomment}
%\newenvironment{reference}{}{}
\newenvironment{slogan}{\comment}{\endcomment}
\newenvironment{history}{\comment}{\endcomment}

% For commutative diagrams we use Xy-pic
\usepackage[all]{xy}

% We use 2cell for 2-commutative diagrams.
\xyoption{2cell}
% Theorem environments.
%
\theoremstyle{plain}
\newtheorem{theorem}[subsection]{Theorem}
\newtheorem{proposition}[subsection]{Proposition}
\newtheorem{lemma}[subsection]{Lemma}

\theoremstyle{definition}
\newtheorem{definition}[subsection]{Definition}
\newtheorem{example}[subsection]{Example}
\newtheorem{exercise}[subsection]{Exercise}
\newtheorem{situation}[subsection]{Situation}

\theoremstyle{remark}
\newtheorem{remark}[subsection]{Remark}
\newtheorem{remarks}[subsection]{Remarks}

\numberwithin{equation}{subsection}

% Macros
%
\def\lim{\mathop{\mathrm{lim}}\nolimits}
\def\colim{\mathop{\mathrm{colim}}\nolimits}
\def\Spec{\mathop{\mathrm{Spec}}}
\def\Hom{\mathop{\mathrm{Hom}}\nolimits}
\def\Ext{\mathop{\mathrm{Ext}}\nolimits}
\def\SheafHom{\mathop{\mathcal{H}\!\mathit{om}}\nolimits}
\def\SheafExt{\mathop{\mathcal{E}\!\mathit{xt}}\nolimits}
\def\Sch{\mathit{Sch}}
\def\Mor{\mathop{\mathrm{Mor}}\nolimits}
\def\Ob{\mathop{\mathrm{Ob}}\nolimits}
\def\Sh{\mathop{\mathit{Sh}}\nolimits}
\def\NL{\mathop{N\!L}\nolimits}
\def\CH{\mathop{\mathrm{CH}}\nolimits}
\def\proetale{{pro\text{-}\acute{e}tale}}
\def\etale{{\acute{e}tale}}
\def\QCoh{\mathit{QCoh}}
\def\Ker{\mathop{\mathrm{Ker}}}
\def\Im{\mathop{\mathrm{Im}}}
\def\Coker{\mathop{\mathrm{Coker}}}
\def\Coim{\mathop{\mathrm{Coim}}}

% Boxtimes
%
\DeclareMathSymbol{\boxtimes}{\mathbin}{AMSa}{"02}

%
% Macros for moduli stacks/spaces
%
\def\QCohstack{\mathcal{QC}\!\mathit{oh}}
\def\Cohstack{\mathcal{C}\!\mathit{oh}}
\def\Spacesstack{\mathcal{S}\!\mathit{paces}}
\def\Quotfunctor{\mathrm{Quot}}
\def\Hilbfunctor{\mathrm{Hilb}}
\def\Curvesstack{\mathcal{C}\!\mathit{urves}}
\def\Polarizedstack{\mathcal{P}\!\mathit{olarized}}
\def\Complexesstack{\mathcal{C}\!\mathit{omplexes}}
% \Pic is the operator that assigns to X its picard group, usage \Pic(X)
% \Picardstack_{X/B} denotes the Picard stack of X over B
% \Picardfunctor_{X/B} denotes the Picard functor of X over B
\def\Pic{\mathop{\mathrm{Pic}}\nolimits}
\def\Picardstack{\mathcal{P}\!\mathit{ic}}
\def\Picardfunctor{\mathrm{Pic}}
\def\Deformationcategory{\mathcal{D}\!\mathit{ef}}
\usepackage{cleveref}
\crefname{originaltheorem}{original theorem}{original theorems}
\Crefname{originaltheorem}{Original theorem}{Original theorems}
\makeatletter
\newlabel{theorem-wedderburn}{{0747}{source}{}{sourcepointer-1}{}}
\newlabel{theorem-wedderburn@cref}{{[originaltheorem][0][]0747}{[1][0][]source}{}{}{}}
\makeatother
\crefname{originallemma}{original lemma}{original lemmas}
\Crefname{originallemma}{Original lemma}{Original lemmas}
\makeatletter
\newlabel{lemma-centralizer}{{0749}{source}{}{sourcepointer-2}{}}
\newlabel{lemma-centralizer@cref}{{[originallemma][0][]0749}{[1][0][]source}{}{}{}}
\makeatother
\title{The center of a finite simple algebra is a finite field extension}
\author{The Stacks Project Authors}
\date{stacks-project:brauer.tex:local-snapshot-20261012; GNU FDL 1.2 or later}
\begin{document}
\maketitle
\noindent\textit{Editorial note.} All delivered mathematical text is the author's own, copied from the pinned Stacks Project chapter file sites-cohomology.tex; no statement or proof step was rewritten, reordered or removed.\par\medskip
\noindent\textit{Editorial note.} The proof is the author's complete argument; both steps name their input and no step is left to the reader.\par\medskip
\noindent\textit{Editorial note.} Every labelled reference occurring inside the counted spans is carried in assembly.external\_references as a non-load-bearing navigation pointer (same-chapter ones at their definition site, cross-chapter ones at the invocation inside the counted span), so no load-bearing mathematics is replaced by a stub. No bibliography citation occurs inside the counted spans.\par\medskip
\noindent\textit{Editorial note.} The same-name check for this card scanned every staged Stacks chapter of this repository -- all sources/stacks-*/pinned-original-source.tex files, the card's own chapter excluded -- and the labels found are recorded in dedup.same\_name\_cross\_chapter on the card.\par\medskip
\noindent\textit{Editorial note.} The preamble is the chapter's own retained minimal preamble (sources/stacks-sites-cohomology/retained-preamble.tex), a verbatim excerpt of the native preamble: the Xy-pic setup and the native theorem declarations with the mathematical macros.\par\medskip
\noindent\textit{Allowed standard prerequisites (editorial).}\begin{itemize}
\item Wedderburn's theorem. Same chapter, Theorem theorem-wedderburn. Supplies the description of the algebra as a matrix algebra over a skew field finite over the ground field.
\item Centralizers commute with tensor products. Same chapter, Lemma lemma-centralizer. Supplies that the center of the matrix algebra is the tensor product of the ground field with the center of the skew field, which identifies the center.
\end{itemize}\medskip
\ifdefined\hypertarget\hypertarget{sourcepointer-1}{}\fi
\noindent\textit{Source reference.} Original theorem 0747: sources/stacks-brauer/pinned-original-source.tex, source lines 358--358. Navigation pointer invoked at proof line 418; the counted span contains the pointer only and the cited result is not reproduced inside it, so the pointer is registered as a non-load-bearing navigation stub in the established form for Stacks material. No bibliography citation is involved.\par\medskip
\ifdefined\hypertarget\hypertarget{sourcepointer-2}{}\fi
\noindent\textit{Source reference.} Original lemma 0749: sources/stacks-brauer/pinned-original-source.tex, source lines 395--395. Navigation pointer invoked at proof line 420; the counted span contains the pointer only and the cited result is not reproduced inside it, so the pointer is registered as a non-load-bearing navigation stub in the established form for Stacks material. No bibliography citation is involved.\par\medskip
\begin{lemma}
\label{lemma-center-csa}
Let $A$ be a finite simple $k$-algebra. Then the center $k'$ of $A$
is a finite field extension of $k$.
\end{lemma}

\begin{proof}
Write $A = \text{Mat}(n \times n, K)$ for some skew field $K$ finite
over $k$, see
Theorem \ref{theorem-wedderburn}.
By
Lemma \ref{lemma-centralizer}
the center of $A$ is $k \otimes_k k'$ where $k' \subset K$ is the
center of $K$. Since the center of a skew field is a field, we win.
\end{proof}

\end{document}
